% theme: quadratic_equations_plane_geometry_applications
\MajorQuestionLesson{平面図形の利用問題}
\SetRowSpace{7mm}

\TypeQuestionSingle{次の問いに答えなさい。}{henka}
\QQ{ある正方形の厚紙がある。この厚紙の縦を2cm、横を4cm長くすると、面積が89cm$^2$の長方形になった。もとの厚紙の1辺の長さを求めなさい。}{}
\QQ{ある正方形の紙がある。この紙の縦を5cm長くし、横を1cm短くすると、面積が112cm$^2$の長方形になった。もとの紙の1辺の長さを求めなさい。}{}
\QQ{縦より横が5cm長い長方形がある。縦を2cm、横を1cm長くすると、面積が140cm$^2$になった。もとの長方形の縦と横の長さを求めなさい。}{}

\TypeQuestionSingle{次の問いに答えなさい。}{shucho}
\QQ{周の長さが40cmで、面積が96cm$^2$の長方形がある。この長方形の縦と横の長さを求めなさい。}{}
\QQ{横の長さが縦の長さより7cm長く、面積が120cm$^2$である長方形がある。この長方形の縦と横の長さを求めなさい。}{}
\QQ{長さ60cmのひもを使って長方形を作ったところ、ひもが8cm余り、長方形の面積は160cm$^2$になった。この長方形の縦と横の長さを求めなさい。}{}

\TypeQuestionSingle{次の問いに答えなさい。}{michi}
\QQ{縦12m、横16mの長方形の土地がある。縦方向と横方向に同じ幅の道を1本ずつ通し、残った部分を花だんにする。花だんの面積が合わせて140m$^2$になるとき、道の幅を求めなさい。}{}
\QQ{縦20m、横30mの長方形の公園の内側に、同じ幅の道を周りに作り、中央を芝生にする。芝生の面積が416m$^2$になるとき、道の幅を求めなさい。}{}
\QQ{縦14m、横18mの長方形の土地がある。縦方向と横方向に同じ幅の道を1本ずつ通し、残った部分を花だんにする。花だんの面積が合わせて140m$^2$になるとき、道の幅を求めなさい。}{}

\LessonPageBreak

\SetRowSpace{7mm}

\TypeQuestionSingle{次の問いに答えなさい。}{kiritori}
\QQ{縦10cm、横16cmの長方形の厚紙がある。四すみから1辺$x$ cmの正方形を切り取ったところ、残った部分の面積は96cm$^2$になった。切り取った正方形の1辺の長さを求めなさい。}{}
\QQ{正方形の花だんの周りに、幅2mの道をつけると、道を含めた全体は1辺12mの正方形になった。もとの花だんの1辺の長さを求めなさい。}{}
\QQ{ある正方形の板から、1辺6cmの正方形を切り取ったところ、残った部分の面積が64cm$^2$になった。もとの正方形の1辺の長さを求めなさい。}{}

\TypeQuestionSingle{次の問いに答えなさい。}{sankaku}
\QQ{底辺が高さより2cm長い三角形があり、面積は40cm$^2$である。この三角形の底辺と高さを求めなさい。}{}
\QQ{上底が$x$ cm、下底が上底より6cm長く、高さが上底と等しい台形がある。この台形の面積が88cm$^2$であるとき、上底、下底、高さを求めなさい。}{}
\QQ{底辺が高さより5cm長い平行四辺形があり、面積は84cm$^2$である。この平行四辺形の底辺と高さを求めなさい。}{}

\TypeQuestionSingle{次の問いに答えなさい。}{kumiawase}
\QQ{1辺$x$ cmの正方形と、縦$x$ cm、横がそれより4cm長い長方形を組み合わせた図形がある。この図形の面積が96cm$^2$であるとき、$x$ の値を求めなさい。}{}
\QQ{1辺の長さが3cm違う2つの正方形がある。2つの正方形の面積の和が65cm$^2$であるとき、それぞれの正方形の1辺の長さを求めなさい。}{}
